% titel: Blatt 10, A9: Titration HCOOH mit NaOH (pH vor/am Äquivalenzpunkt)
% skript: Mortimer Kap. 17 (Titration, Puffer, Hydrolyse)

% Teil I: Werkzeuge, kein Lösungsweg.

Beim Titrieren einer schwachen Säure mit einer starken Base ändert sich der
pH je nach Stadium: vor dem Äquivalenzpunkt wirkt das Gemisch als Puffer,
am Äquivalenzpunkt liegt die konjugierte Base (hydrolysiert) vor.

\begin{werkzeugkasten}[Titration einer schwachen Säure]
\begin{itemize}
  \item \textbf{Vor dem Äquivalenzpunkt:} Es bleiben HA und A$^-$
        nebeneinander $\Rightarrow$ Puffer:
        \[
          pH = pK_s + \lg\frac{n(\mathrm{A^-})}{n(\mathrm{HA})},
        \]
        mit $n(\mathrm{A^-})$ aus der zugegebenen Base.
  \item \textbf{Am Äquivalenzpunkt:} alles HA ist neutralisiert, es liegt
        nur die schwache Base A$^-$ vor:
        \[
          pOH \approx \frac12(pK_b - \lg c)\,,\quad
          pK_b = 14-pK_s,\qquad [OH^-] = \sqrt{K_b c}.
        \]
  \item Stöchiometrie $n(\mathrm{OH^-})=n(\mathrm{HA})$ bestimmt die
        benötigte NaOH-Volumen.
\end{itemize}
\textit{Skript:} Mortimer Kap. 17 (Titrationskurven).
\end{werkzeugkasten}

\begin{achtung}[Darauf achten]
\begin{itemize}
  \item Am Äquivalenzpunkt ist der pH \emph{nicht} 7: Die konjugierte Base
        (Formiat) macht die Lösung basisch.
  \item pK$_s$(HCOOH)$=3.75 \Rightarrow pK_b=10.25$.
  \item Für Teil a) erst die im Überschuss verbleibende Säuremenge und das
        gebildete A$^-$ bestimmen (Stoffmengen, nicht Konzentrationen nötig
        für das Ratio).
  \item Gesamtvolumen ändert sich; für pH/Puffer kürzt es sich im
        Verhältnis häufig, für $c$ am Äquivalenzpunkt aber nicht.
\end{itemize}
\end{achtung}
